% BHU PYQ -- Transcribed & Corrected
% Paper: CSCMJ45 / CSCMN41: Numerical Computing
% Exam: B.Sc. (Hons.) Semester IV Examination, 2025-26 (NEP)
% Category: Major & Minor Course
% Department: Computer Science

\documentclass[11pt,a4paper]{article}
\usepackage[a4paper,top=2.5cm,bottom=2.5cm,left=2.8cm,right=2.8cm,headheight=14pt]{geometry}
\usepackage{amsmath,amssymb}
\usepackage[T1]{fontenc}
\usepackage[utf8]{inputenc}
\usepackage{lmodern,microtype,fancyhdr,enumitem,hyperref,lastpage,parskip,booktabs}

\hypersetup{
    colorlinks=true,
    urlcolor=black,
    linkcolor=black,
    pdftitle={CSCMJ45 / CSCMN41: Numerical Computing -- BHU Sem IV 2025-26 (NEP)}
}

\pagestyle{fancy}
\fancyhf{}
\renewcommand{\headrulewidth}{0.4pt}
\renewcommand{\footrulewidth}{0.4pt}
\fancyhead[L]{\small   \textit{Banaras Hindu University}}
\fancyhead[R]{\small   \textit{CSCMJ45/CSCMN41: Numerical Computing (Sem IV 2025-26)}}
\fancyfoot[C]{\small Page  \thepage\ of \pageref{LastPage}}


\newlist{parts}{enumerate}{2}
\setlist[parts,1]{label=(\alph*),leftmargin=*,itemsep=4pt,topsep=2pt}
\setlist[parts,2]{label=(\roman*),leftmargin=*,itemsep=2pt,topsep=2pt}

\newcommand{\pts}[1]{\hfill   \textit{[#1]}}


\begin{document}


\begin{center}
  {\large\bfseries BANARAS HINDU UNIVERSITY}\\[2pt]
  {\normalsize B.Sc. (Hons.) --- Semester IV Examination, 2025-26 (NEP)}\\[6pt]
  \rule{0.8\linewidth}{0.5pt}\\[6pt]
  {\large\bfseries COMPUTER SCIENCE}\\[4pt]
  {\large\bfseries Paper Code: CSCMJ45 / CSCMN41 : Numerical Computing}\\[4pt]
  {\normalsize Time: 2:30 Hours\hfill Maximum Marks: 70}
\end{center}

\rule{\linewidth}{0.4pt}

\smallskip



\noindent   \textit{Instructions: Attempt any   \textbf{5} questions.}

\bigskip




\noindent   \textbf{Question 1.} Derive the general term for evaluation and hence write the algorithm for:\pts{14}

\begin{parts}
  \item Polynomial Evaluation
  \item Polynomial Division
\end{parts}

\medskip\rule{\linewidth}{0.2pt}\medskip




\noindent   \textbf{Question 2.} Find a real root of the equation $x^3 + x^2 + x + 7 = 0$, correct to 4 decimal places using:\pts{7 + 7}

\begin{parts}
  \item Regula-Falsi method
  \item Newton-Raphson method
\end{parts}



\noindent Compare the number of steps taken in each case.

\medskip\rule{\linewidth}{0.2pt}\medskip




\noindent   \textbf{Question 3.} Solve the following system of linear equations by (a) Gauss-elimination and (b) Gauss-Jordan methods:\pts{7 + 7}
\[
  
\begin{aligned}
    9x_1 + 2x_2 + 4x_3 &= 20 \\
    x_1 + 10x_2 + 4x_3 &= 6 \\
    2x_1 - 4x_2 + 10x_3 &= -15
  \end{aligned}
\]

\medskip\rule{\linewidth}{0.2pt}\medskip




\noindent   \textbf{Question 4.} Decompose the matrix
\[
  A = 
\begin{bmatrix} 5 & 2 & -1 \\ 7 & 1 & -5 \\ 3 & 7 & 4 \end{bmatrix}
\]
into the form $LU$ and hence solve the system $AX = b$ where $b = [4 \quad 8 \quad 10]^T$.\pts{7 + 7}

\medskip\rule{\linewidth}{0.2pt}\medskip




\noindent   \textbf{Question 5.} From the table given below, calculate $f(1.023)$ and $f(1.235)$ using Newton-Gregory forward and Lagrange interpolation formula respectively:\pts{7 + 7}

\begin{center}
  
\begin{tabular}{@{}ccccccc@{}}
     oprule
    $x$ & $1.00$ & $1.05$ & $1.10$ & $1.15$ & $1.20$ & $1.25$ \\
    \midrule
    $f(x)$ & $0.682689$ & $0.706282$ & $0.728668$ & $0.749586$ & $0.769861$ & $0.788700$ \\
    ottomrule
  \end{tabular}
\end{center}

\medskip\rule{\linewidth}{0.2pt}\medskip




\noindent   \textbf{Question 6.} Solve the given differential equation $\frac{dy}{dx} = x^2 + y$ with $y(0) = 1$, to compute $y(0.25)$ using step size $h = 0.05$:\pts{7 + 7}

\begin{parts}
  \item Euler's method
  \item Heun's method
\end{parts}



\noindent Compare the results obtained in each case.

\medskip\rule{\linewidth}{0.2pt}\medskip




\noindent   \textbf{Question 7.} Answer the following:\pts{7 + 7}

\begin{parts}
  \item Calculate the following integral using Simpson's $1/3$ rule with $h = 0.25$:
  \[
    \int_{0}^{2} \frac{1}{x^3 + x + 1} \, dx
  \]
  \item Integrate $f(x) = 2x^3 - 3x^2 + 4x - 5$ from $-2$ to $4$ using Gauss-Legendre 2-point formula.
\end{parts}

\vfill
\rule{\linewidth}{0.4pt}

\begin{center}\small
\href{https://bhu-exam.vercel.app/}{Click here to visit our website}\\[4pt]
\href{https://me-aryan.vercel.app/}{Click here to visit our contributor}
\end{center}

\end{document}
